\documentclass[12pt]{article}

\usepackage[margin=1in]{geometry}
\usepackage{mathptmx}
\usepackage{setspace}
\usepackage{lineno}
\usepackage{microtype}
\usepackage{amsmath,amssymb}
\usepackage{booktabs}
\usepackage{graphicx}
\usepackage[super,sort&compress]{natbib}
\usepackage[hidelinks]{hyperref}
\usepackage{xcolor}

\doublespacing
\linenumbers

% Draft-only macros. Remove before submission.
\newcommand{\result}[1]{\textbf{\textcolor{blue}{[RESULT: #1]}}}
\newcommand{\methodtodo}[1]{\textbf{\textcolor{purple}{[METHOD: #1]}}}
\newcommand{\datatodo}[1]{\textbf{\textcolor{teal}{[DATA: #1]}}}
\newcommand{\paneltodo}[2]{%
  \fbox{%
    \begin{minipage}[t]{#1}
    \raggedright\small
    #2
    \end{minipage}%
  }%
}

\newcommand{\figureonetodo}{%
\begin{figure}[p]
\centering
\paneltodo{0.56\textwidth}{%
\textbf{A | Spectrum as partial observation}\\[3pt]
\textbullet~Precursor + ionization / collision condition\\
\textbullet~Latent multi-branch fragmentation World\\
\textbullet~Only part of the World is observed as MS/MS\\[2pt]
\centering $M,c \rightarrow \mathcal{W}(M,c) \rightarrow S_{\rm obs}$
}
\hfill
\paneltodo{0.38\textwidth}{%
\textbf{B | One World, three queries}\\[3pt]
\textbullet~Free: $p(\mathcal T,S\mid M,c)$\\
\textbullet~Guided: $p(\mathcal T\mid M,S_{\rm obs},c)$\\
\textbullet~Inverse: $p(M,\mathcal T\mid S_{\rm obs},c)$\\
\textbullet~All share one chemical state / transition system
}

\vspace{5pt}

\paneltodo{0.30\textwidth}{%
\textbf{C | Autoregressive 1e chemistry}\\[3pt]
\textbullet~source $\rightarrow$ target $\rightarrow$ 1e\\
\textbullet~STOP / REPEAT / NEW\\
\textbullet~2 repeated 1e events recover a 2e transfer\\
\textbullet~No fragmentation-template catalogue
}
\hfill
\paneltodo{0.34\textwidth}{%
\textbf{D | Probability-mass semantics}\\[3pt]
\textbullet~valid chemistry\\
\textbullet~invalid chemistry\\
\textbullet~censored / unsampled mass\\
\textbullet~learned STOP\\[2pt]
\centering $P_{\rm valid}+P_{\rm invalid}+P_{\rm censored}+P_{\rm unsampled}+P_{\rm stop}=1$
}
\hfill
\paneltodo{0.30\textwidth}{%
\textbf{E | Multi-branch reaction DAG}\\[3pt]
\textbullet~Sibling fragmentation branches\\
\textbullet~Branch-local STOP\\
\textbullet~Reconvergent canonical states\\
\textbullet~Conditional and path probability
}

\vspace{5pt}

\paneltodo{0.94\textwidth}{%
\textbf{F | Outer-loop scientific evolution}\\[3pt]
failure $\rightarrow$ bounded proposal $\rightarrow$ typed chemistry compiler / conservation checks $\rightarrow$ molecule-disjoint held-out gate $\rightarrow$ promoted World\\
\textit{Keep this as a thin supporting strip; the agent is outside online chemistry inference.}
}
\caption{\textbf{Figure 1 layout placeholder | A molecular fragmentation world unifies prediction, explanation and inverse inference.}
Top: latent World and W/G/V queries. Middle: electron-event language, probability semantics and reaction DAG. Bottom: supporting outer-loop evolution.}
\end{figure}
}

\newcommand{\figuretwotodo}{%
\begin{figure}[p]
\centering
\paneltodo{0.45\textwidth}{%
\textbf{A | Free World performance}\\[3pt]
\textbullet~Reaction-network example\\
\textbullet~Observed mass coverage\\
\textbullet~Intensity-weighted coverage\\
\textbullet~Ordered-transition / multistage coverage
}
\hfill
\paneltodo{0.45\textwidth}{%
\textbf{B | Guided vs Free under matched compute}\\[3pt]
\textbullet~Same molecule / chemistry / budget\\
\textbullet~Only Guided observes the spectrum\\
\textbullet~Coverage vs expanded states\\
\textbullet~States or wall time to matched coverage
}

\vspace{5pt}

\paneltodo{0.38\textwidth}{%
\textbf{C | Inverse posterior structures}\\[3pt]
\textbullet~Spectrum $\rightarrow$ multiple candidates\\
\textbullet~Exact / Recall@K\\
\textbullet~Best similarity / MCES\\
\textbullet~Validity + posterior diversity
}
\hfill
\paneltodo{0.52\textwidth}{%
\textbf{D | Bidirectional consistency}\\[3pt]
\centering $S_{\rm obs}\rightarrow M\rightarrow\mathcal W(M)\rightarrow\hat S$\\[3pt]
\raggedright
\textbullet~GT vs same-formula / near-isomer decoys\\
\textbullet~Executable path and intensity support\\
\textbullet~Ordered-transition consistency\\
\textbullet~Unresolved World probability mass
}

\vspace{5pt}

\paneltodo{0.58\textwidth}{%
\textbf{E | Graded structural evidence}\\[3pt]
formula $\rightarrow$ valid graph $\rightarrow$ World reachable $\rightarrow$ path support\\
$\rightarrow$ candidate-specific evidence $\rightarrow$ bidirectional consistency\\
$\rightarrow$ calibrated bounded statement $\rightarrow$ authentic standard\\[3pt]
\textbullet~Overlay decreasing coverage and increasing reliability
}
\hfill
\paneltodo{0.32\textwidth}{%
\textbf{F | Prospective credibility}\\[3pt]
\textbullet~Historical freeze date\\
\textbullet~Dark at freeze, identified later\\
\textbullet~Containment / calibration\\
\textbullet~Structural-distance shift
}
\caption{\textbf{Figure 2 layout placeholder | World-based inference establishes graded structural evidence.}
Top: Free and Guided capability. Middle: inverse posterior and bidirectional verification. Bottom: evidence hierarchy and prospective calibration.}
\end{figure}
}

\newcommand{\figurethreetodo}{%
\begin{figure}[p]
\centering
\paneltodo{0.94\textwidth}{%
\textbf{A | Repository-scale evidence funnel}\\[3pt]
all dark spectra $\rightarrow$ formula compatible $\rightarrow$ World reachable $\rightarrow$ executable path support\\
$\rightarrow$ candidate-specific evidence $\rightarrow$ bidirectional consistency $\rightarrow$ calibrated bounded structure\\[3pt]
\textbullet~Show denominator at every step\\
\textbullet~Keep unresolved and computationally censored branches visible
}

\vspace{5pt}

\paneltodo{0.45\textwidth}{%
\textbf{B | Spectral observations to structural entities}\\[3pt]
\textbullet~Adduct / isotope / charge-state collapse\\
\textbullet~In-source fragment / redundant acquisition collapse\\
\textbullet~Over- and under-merging validation\\
\textbullet~Entity-count uncertainty
}
\hfill
\paneltodo{0.45\textwidth}{%
\textbf{C | Recurrence across independent datasets}\\[3pt]
\textbullet~Independent datasets / laboratories only\\
\textbullet~Entity / family recurrence distribution\\
\textbullet~Separate single-study and recurrent chemistry
}

\vspace{5pt}

\paneltodo{0.45\textwidth}{%
\textbf{D | Evidence strength vs recurrence}\\[3pt]
\textbullet~Recurrence category on x-axis\\
\textbullet~Evidence level / reliability on y-axis\\
\textbullet~Test whether recurrent chemistry is more resolvable
}
\hfill
\paneltodo{0.45\textwidth}{%
\textbf{E | Near-known vs remote chemistry}\\[3pt]
\textbullet~Frozen characterized-metabolite reference\\
\textbullet~Distance-to-known distribution\\
\textbullet~Near-known halo / remote families / insufficient resolution
}

\vspace{5pt}

\paneltodo{0.94\textwidth}{%
\textbf{F | Representative high-evidence dark chemistry}\\[3pt]
3--5 examples only: World-resolved structure | bounded isomer set | remote recurrent family | standard-confirmed anchor\\
For each example show: structure + evidence level + recurrence + remaining uncertainty.
}
\caption{\textbf{Figure 3 layout placeholder | High-evidence inference resolves the dark metabolome at scale.}
The full-width evidence funnel is the dominant panel; lower panels define entities, recurrence, novelty and representative chemistry.}
\end{figure}
}

\newcommand{\figurefourtodo}{%
\begin{figure}[p]
\centering
\paneltodo{0.58\textwidth}{%
\textbf{A | Structural landscape}\\[3pt]
\textbullet~Known metabolites + high-evidence dark entities\\
\textbullet~Structural relationships, not spectral embedding alone\\
\textbullet~Near-known halo and remote recurrent islands\\
\textbullet~Use family boundaries / density, not a network hairball
}
\hfill
\paneltodo{0.32\textwidth}{%
\textbf{B | Recurrent family architecture}\\[3pt]
\textbullet~Family-size distribution\\
\textbullet~Recurrence distribution\\
\textbullet~Near-known vs remote composition
}

\vspace{5pt}

\paneltodo{0.94\textwidth}{%
\textbf{C | Public perturbation evidence}\\[3pt]
germ-free / gnotobiotic | antibiotic | defined diet\\
Show manually verified study count, biological sample count and platform for each branch.
}

\vspace{5pt}

\paneltodo{0.58\textwidth}{%
\textbf{D | Source dependence over chemical space}\\[3pt]
\textbullet~Reuse exactly the coordinates from panel A\\
\textbullet~Overlay microbiota-dependent / diet-dependent / host-associated\\
\textbullet~Keep mixed and unresolved families visible
}
\hfill
\paneltodo{0.32\textwidth}{%
\textbf{E | Structure--source coupling}\\[3pt]
\textbullet~Local label autocorrelation\\
\textbullet~Motif / class enrichment\\
\textbullet~Cross-study consistency\\
\textbullet~Permutation null
}

\vspace{5pt}

\paneltodo{0.94\textwidth}{%
\textbf{F | Worked remote family}\\[3pt]
family structures + recurrence across datasets + perturbation effect + source-dependence label + evidence level\\
Use dependent / associated language; do not imply biosynthesis without direct evidence.
}
\caption{\textbf{Figure 4 layout placeholder | Resolved dark chemistry reveals structural and biological organization.}
Panels A and D use the same chemical-space coordinates so biological source can be read directly against structural organization.}
\end{figure}
}

\newcommand{\figurefivetodo}{%
\begin{figure}[p]
\centering
\paneltodo{0.47\textwidth}{%
\textbf{A | Global feature-versus-family test}\\[3pt]
\textbullet~Same samples / preprocessing / covariates\\
\textbullet~Anonymous features vs frozen structural families\\
\textbullet~Primary endpoint: replication or sign concordance\\
\textbullet~Show paired effect with confidence interval
}
\hfill
\paneltodo{0.43\textwidth}{%
\textbf{B | Reproducibility gain}\\[3pt]
\textbullet~Replication\\
\textbullet~Effect-sign concordance\\
\textbullet~Multiplicity-controlled yield\\
\textbullet~Effect stability across all eligible contrasts
}

\vspace{5pt}

\paneltodo{0.39\textwidth}{%
\textbf{C | Principal biological programme}\\[3pt]
\textbullet~Discovery cohort\\
\textbullet~Independent replication\\
\textbullet~Sample-level effects\\
\textbullet~Adjusted significance / effect CI
}
\hfill
\paneltodo{0.51\textwidth}{%
\textbf{D | Family chemistry}\\[3pt]
\textbullet~Molecular structures / bounded statements\\
\textbullet~Evidence level + confidence + recurrence\\
\textbullet~Interpretable structural differences\\
\textbullet~Do not draw a pathway without biosynthetic evidence
}

\vspace{5pt}

\paneltodo{0.45\textwidth}{%
\textbf{E | Authentic-standard anchors}\\[3pt]
\textbullet~Sample vs standard MS/MS overlay\\
\textbullet~Retention / coelution\\
\textbullet~Precursor agreement\\
\textbullet~Clearly mark standard-confirmed evidence
}
\hfill
\paneltodo{0.45\textwidth}{%
\textbf{F | Explicitly unresolved member}\\[3pt]
\textbullet~Surviving isomer set\\
\textbullet~Evidence eliminating alternatives\\
\textbullet~Why current MS/MS cannot separate survivors\\
\textbullet~Predicted next discriminating measurement
}
\caption{\textbf{Figure 5 layout placeholder | Dark molecular families reveal biology hidden from feature-level metabolomics.}
The global paired analysis comes first; one replicated programme then connects biological effect, molecular family, orthogonal standards and residual uncertainty.}
\end{figure}
}

\title{\textbf{A molecular fragmentation world reveals the dark metabolome}}
\author{Author names and affiliations to be finalized}
\date{}

\begin{document}
\maketitle

\noindent\textbf{
Tandem mass spectra are incomplete observations of the gas-phase reaction processes that generate them, yet molecular structure elucidation is usually learned as a direct mapping between spectra and molecular candidates. Here we instead learn a probabilistic molecular fragmentation world in which ionization, electron-flow events, hydrogen transfer, competing fragmentation branches and their observations share one conservation-preserving state space. The same World supports three complementary modes of inference: Free prediction of molecular fragmentation, Guided search for reaction paths that explain an observed spectrum, and Inverse generation of molecular structures that must remain consistent when replayed through the forward World. Under molecule-disjoint evaluation, the learned World recovers \result{Free mass/intensity/transition result}; Guided inference improves \result{matched-budget explanatory endpoint} by \result{Guided gain}; and Inverse inference recovers the true structure within the top \result{K} candidates in \result{Inverse Recall@K}, while maintaining \result{posterior diversity/validity}. Bidirectional spectrum--structure consistency converts molecular prediction into a hierarchy of auditable evidence, and stronger evidence levels correspond to \result{reliability gradient}, including \result{prospective containment/calibration} on structures deposited only after the system was frozen. Applying this evidence hierarchy to \result{dark spectrum count} previously unannotated tandem mass spectra resolves \result{evidence-level distribution/high-evidence count} into bounded structural statements and \result{recurrent entity/family count} recurrent molecular entities or families. These resolved dark structures form \result{headline structural-organization result}, partition reproducibly by \result{microbiota/diet/host source result}, and improve \result{feature-versus-family biological reproducibility endpoint} relative to feature-level analysis, revealing \result{principal biological programme}. Our results establish molecular fragmentation as a learnable world for forward and inverse chemical inference and turn the dark metabolome from a collection of unidentified spectra into an experimentally bounded and biologically analysable chemical space.
}

\vspace{1em}

Mass spectrometry is the principal technology for surveying small molecules in biology, medicine and the environment.\cite{Wang2016GNPS,Alseekh2021,Gentry2024} Yet most tandem mass spectra generated by untargeted studies remain without a confident structural assignment.\cite{daSilva2015,Alseekh2021} Recent advances in spectral retrieval, molecular networking, formula and fingerprint prediction, learned spectrum models and de novo structure generation have substantially improved the hypotheses that can be proposed for a spectrum.\cite{Duhrkop2015,Duhrkop2019,Duhrkop2021,Wang2021CFMID,Goldman2023,Bohde2025,Qiang2026DeepMet} The dominant abstraction, however, remains a mapping between an observed spectrum and an answer: a score, a candidate rank, a fingerprint, or a molecular structure.

A tandem mass spectrum is not the reaction process itself. It is a partial observation of an ensemble of gas-phase events that depend on molecular structure, ionization state, collision conditions, competing pathways and instrument response. A fragment absent from an acquired spectrum is therefore not automatically chemically impossible, just as a peak present in the spectrum does not by itself specify the unique trajectory that produced it. Treating the spectrum as the complete target obscures this latent reaction process and makes forward prediction, mechanistic explanation and inverse structure inference separate problems that need not agree with one another.

We take a different view. For a molecule \(M\) under experimental condition \(c\), we define a latent molecular fragmentation world \(\mathcal{W}(M,c)\): a probabilistic, multi-branch reaction network over ionization microstates, electron-flow events, charged and neutral products, hydrogen transfer, branch competition and termination. The observed spectrum \(S_{\mathrm{obs}}\) is generated through an observation process acting on this world,
\[
M,c \rightarrow \mathcal{W}(M,c) \rightarrow S_{\mathrm{obs}}.
\]
This formulation separates chemical possibility from experimental visibility and makes the latent reaction world, rather than a vector of spectral bins, the shared object of inference.

ORBIT-MS implements this idea as a World--Generator--Verifier system. The World learns Free fragmentation, \(p(\mathcal{T},S\mid M,c)\), without access to a target spectrum. The Verifier performs Guided inference, \(p(\mathcal{T}\mid M,S_{\mathrm{obs}},c)\), allocating finite computation toward branches that explain or discriminate the observed evidence. The Generator performs Inverse inference, \(p(M,\mathcal{T}\mid S_{\mathrm{obs}},c)\), maintaining multiple molecular hypotheses that must remain chemically valid and are replayed through the same forward World. These are not three unrelated scorers: they are three conditional queries to one chemical process.

The combinatorial size of this world makes reinforcement learning a natural inference mechanism rather than an end in itself. Free fragmentation requires probability flow over many valid trajectories; Guided explanation requires long-horizon allocation of a finite expansion budget; Inverse inference requires a diverse posterior over molecular structures rather than a single winner. GFlowNet, offline value learning, data aggregation and sequential Monte Carlo therefore determine where computation is spent, while chemistry determines which actions are legal. Atom, electron, charge, hydrogen, exact-mass, formula and valence constraints are enforced by the executor rather than traded against reward.

This shared World also changes the form of structural evidence. A candidate can progress from mass or formula compatibility to chemical validity, World reachability, executable path support, candidate-specific discriminative evidence and bidirectional spectrum--structure consistency. Because stronger evidence can be evaluated against known structures and calibrated prospectively, the output need not collapse every spectrum to an unjustified rank-1 identity. Instead, the system reports the strongest bounded structural statement supported by the available evidence.

We use this capability to address the dark metabolome. Public repositories contain enormous numbers of unannotated spectra,\cite{Wang2016GNPS,Bushuiev2026Dreams} but their chemistry is usually analysed as anonymous features or spectral clusters. Once dark spectra can be moved to explicit evidence levels and bounded molecular structures, they can instead be counted as recurrent chemical entities, organized into structural families, related to characterized metabolism, assigned perturbation-dependent biological context and compared as molecular objects across cohorts. The paper therefore proceeds from a fragmentation-world paradigm, through graded structural evidence, to large-scale analysis of dark chemistry and its biology.

\section*{A molecular fragmentation world unifies prediction, explanation and inverse inference}

We first tested the central premise that the latent object behind a tandem mass spectrum can be learned as a probabilistic reaction world rather than only as a direct structure-to-spectrum mapping. The World represents ionization microstates, explicit atom and hydrogen identities, bond-electron states, formal charge, radicals, exact mass, charged and neutral coproducts and branch-local termination. Electron-flow actions are executed only when they satisfy hard conservation and materialization constraints; learned policies therefore estimate preference and probability over legal chemistry rather than learning legality from penalties.

On \result{World held-out cohort and denominator} molecule-disjoint held-out molecules, the Free World generated multi-branch fragmentation networks that covered \result{Free observed-mass coverage} of observed fragment masses, \result{Free intensity-weighted coverage} of experimental fragment intensity and \result{Free ordered-transition coverage} of observed parent--child transitions, while maintaining a maximum probability-conservation error of \result{probability-conservation error} \figureonetodo. Multistage expansion contributed \result{multistage-added coverage} evidence not reachable by one-step fragmentation. In contrast to a direct spectrum predictor, these outputs retain the reaction states and trajectories that connect molecular structure to each supported observation \result{direct-predictor comparison}.

The multi-branch representation was essential. Under the same transition model and matched expansion budget, retaining sibling branches, branch-local STOP events and reconvergent paths improved \result{multibranch endpoint} from \result{single-path result} to \result{multi-branch result}. The remaining probability on unexpanded frontier states was reported as censored mass rather than interpreted as chemical impossibility. This distinction allowed the World to represent multiple competing fragmentation routes without forcing a single mechanistic explanation.

The same state and transition semantics were then queried in three directions. Free inference predicted the reaction world from a molecule without reading the target spectrum. Guided inference conditioned on an observed spectrum and learned which frontier states and electron-flow events were worth expanding under a fixed computational budget. Inverse inference constructed multiple candidate molecules from spectral and precursor constraints and required completed candidates to remain consistent when passed forward through the World. A shared-World ablation showed \result{shared-world coupling result}, indicating that the coupling is not simply architectural parameter sharing but a common chemical state and transition system.

The distinction between chemical legality and learned preference was enforced throughout. Across \result{number of challenged transitions} deliberately legal and invalid transitions, atom, electron, charge, hydrogen and exact-mass constraints rejected \result{invalid rejection rate} of invalid actions while preserving \result{legal execution rate} of valid actions. Forward/reverse closed-system tests recovered the original state in \result{round-trip rate} of cases for transitions with a defined inverse. These checks establish the executable boundary of the World; they do not imply that every allowed event is a validated gas-phase mechanism, whose propensity remains learned and empirically evaluated.

Together, these results define one chemical World with three modes of inference: prediction, explanation and inverse structure inference.

\section*{World-based inference establishes graded structural evidence}

We next asked what the shared World adds to molecular structure inference. A direct generator can propose a molecule, but the proposal does not by itself state why that molecule is compatible with the observation or which alternatives the spectrum can exclude. WGV instead evaluates a sequence of increasingly specific evidence: mass/formula compatibility, chemical validity, World reachability, executable path support, candidate-specific discrimination, bidirectional spectrum--structure consistency and finally a calibrated bounded structural statement.

The Free World provides the first mechanistic layer. On the held-out cohort, it recovered \result{Free evidence result} while preserving the target-spectrum firewall. Guided inference then received the same molecule and the observed spectrum but no additional chemistry. Under an identical budget of \result{matched compute budget}, Guided increased \result{primary Guided endpoint} from \result{Free matched-budget result} to \result{Guided matched-budget result}, and required \result{compute reduction} fewer expanded states to reach matched \result{coverage target} \figuretwotodo. The gain arose from reallocating computation toward evidence-bearing branches rather than from changing which chemical transitions were legal.

Inverse inference addressed the complementary problem. On \result{Inverse held-out cohort} molecule-disjoint spectra, the Generator produced chemically valid and formula-consistent candidates in \result{validity rate} of cases, recovered the true structure at Exact@1/@5/@10 of \result{Exact@1}, \result{Exact@5} and \result{Exact@10}, and achieved a mean best structural similarity of \result{best Tanimoto/MCES}. The posterior retained \result{diversity result} distinct structural modes rather than collapsing immediately to one candidate. Removing World guidance or forward replay changed \result{inverse ablation result}, establishing that the inverse model benefits from inference through the learned reaction process rather than only from graph validity constraints.

We then evaluated completed candidates by replaying them through the same World. In same-formula and near-isomer pools, the ground-truth candidate exceeded the strongest decoy in \result{bidirectional discrimination endpoint} of cases when scored by forward path/intensity consistency, compared with \result{mass/similarity comparator result} for weaker evidence. The advantage was concentrated in spectra containing candidate-specific paths that could be executed for the true structure but not for its closest alternatives. This closed loop,
\[
S_{\mathrm{obs}}\rightarrow M\rightarrow \mathcal{W}(M)\rightarrow \hat S,
\]
turns a molecular proposal into a falsifiable hypothesis about how the observation could have arisen.

Structural reliability increased with evidence strength. Across known-structure evaluation sets, true-structure containment rose from \result{reliability at low evidence} at mass/formula compatibility to \result{reliability at World reachability}, \result{reliability at candidate-specific evidence} and \result{reliability at bidirectional consistency} at successively stronger levels, with corresponding changes in candidate-set size and resolution. We froze these evidence definitions before repository-scale deployment.

Finally, we tested whether the highest model-derived evidence remained honest outside the chemistry available at development time. The complete WGV system and calibration mapping were frozen against a \datatodo{historical library freeze date} snapshot, then evaluated on \result{prospective cohort size} spectra that were dark at freeze but whose structures were independently deposited later. At \result{stated confidence}, the true structure was contained in the reported bounded statement in \result{prospective containment}, with calibration error \result{prospective calibration error}; calibration under increasing structural distance changed by \result{distance-shift calibration result}. Calibrated comparator outputs on the same set showed \result{comparator calibration result}. These results justify using evidence level, rather than rank alone, as the unit of structural confidence in the dark-metabolome analysis.

\section*{High-evidence inference resolves the dark metabolome at scale}

We applied the frozen WGV system and evidence hierarchy to \result{dark spectrum count} tandem mass spectra that lacked an accepted spectral-library match under the frozen search protocol. Sample, tissue, perturbation and phenotype labels were unavailable to structural inference. Each spectrum was assigned the strongest evidence level it reached, including unresolved outcomes and computational censoring.

The resulting evidence funnel showed that \result{formula-compatible count/fraction} spectra reached mass/formula compatibility, \result{World-reachable count/fraction} were supported by the learned World, \result{path-supported count/fraction} had executable paths explaining observed evidence, \result{candidate-specific count/fraction} reached candidate-specific discrimination, and \result{bidirectionally consistent count/fraction} satisfied spectrum--structure--World consistency \figurethreetodo. At the frozen calibrated operating point, \result{bounded-statement count/fraction} received a bounded structural statement at \result{evidence-level distribution}, while the remainder stayed at weaker levels rather than being forced to a unique structure.

We then asked how many chemical entities these observations represent. Entity-collapsing rules for adducts, isotopologues, charge states, in-source fragments and redundant acquisitions were frozen and validated on known compounds, yielding over- and under-merging rates of \result{collapse validation rates}. Propagating these errors, \result{dark observation count} dark observations corresponded to \result{high-evidence structural entity count with interval} high-evidence structural entities, a \result{feature-to-entity contraction} contraction relative to the raw observation count. This count is secondary to the evidence funnel: it describes only chemistry that has reached a stated structural evidence level.

A substantial subset was reproducible across independent datasets. \result{recurrent entity/family count} entities or families occurred in at least \result{recurrence threshold} independent datasets, and \result{high-recurrence count} recurred across \result{higher recurrence threshold} or more. Recurrence was associated with \result{evidence-recurrence relationship}, suggesting that reproducible dark chemistry is not simply the least structurally informative portion of the corpus.

We next positioned the recurrent high-evidence population relative to characterized metabolism using a reference space and structural-distance definition frozen before counting. \result{near-known fraction} fell within the predefined near-known halo, whereas \result{remote recurrent fraction} formed recurrent structural families outside that neighbourhood; \result{insufficient-resolution fraction} could not be assigned a precise structural distance and remained separate. Thus the high-evidence dark population contains both chemistry adjacent to existing metabolite knowledge and recurrent structural regimes poorly represented by current references.

\section*{Resolved dark chemistry reveals structural and biological organization}

The recurrent high-evidence population was not a random set of isolated structures. A structural relationship graph over characterized metabolites and resolved dark entities revealed \result{structural-landscape result}, including \result{remote family count} recurrent remote families with \result{family-size/recurrence result} \figurefourtodo. These families were defined from chemistry alone, before any biological labels were examined.

We then asked whether biological source organizes this structural landscape. We assembled \datatodo{number/accessions of manually verified perturbation datasets} public mouse studies containing germ-free or gnotobiotic comparisons, antibiotic perturbation or defined dietary intervention. The attribution pipeline first recovered expected effects for known source-dependent metabolites in \result{positive-control recovery}, while label permutation yielded \result{permutation-null result}. Only after this gate were dark families interpreted.

Among \result{eligible recurrent-family denominator} eligible families, \result{microbiota-dependent fraction} showed reproducible microbiota dependence, \result{diet-dependent fraction} diet dependence and \result{host-associated fraction} host association; \result{mixed fraction} were mixed and \result{source-unresolved fraction} remained unresolved. These labels denote perturbation dependence rather than direct biosynthetic origin.

Source labels were non-randomly distributed over structural space. Related families shared source dependence more often than expected under permutation by \result{structure-source coupling statistic}; microbiota-dependent regions were enriched for \result{microbiota motifs/classes}, whereas diet-dependent regions were enriched for \result{diet motifs/classes}. The effect remained \result{cross-study/leave-one-study-out result} across independent datasets. A representative \result{worked remote family} family combined \result{family structural core}, recurrence across \result{family recurrence} datasets and a consistent \result{source perturbation effect}, illustrating how structural resolution converts anonymous perturbation-associated peaks into a chemically coherent dark family.

These results indicate that the newly resolved dark chemical population has both structural and biological organization.

\section*{Dark molecular families reveal biology hidden from feature-level metabolomics}

Finally, we asked whether the structural information supplied by WGV changes biological inference. We froze structural-family membership and all eligible biological contrasts before comparing two analyses of the same samples: one treating each dark feature as an independent analyte and one treating chemically related observations as frozen molecular families.

Across \result{number of eligible cohorts/contrasts}, family-level analysis changed the preregistered primary reproducibility endpoint from \result{feature-level primary endpoint} to \result{family-level primary endpoint}, a paired gain of \result{global family gain with confidence interval} \figurefivetodo. Secondary analyses showed \result{effect-sign concordance result}, \result{replication-rate result} and \result{multiplicity-controlled association-yield result}. Because samples, covariates and preprocessing were identical between the two arms, this comparison isolates the effect of changing the analytical object from anonymous features to high-evidence structural families.

The strongest frozen programme, \result{principal family}, was associated with \result{phenotype/contrast} in \result{discovery result} and reproduced in \result{independent replication result}. Its \result{family member count} members share \result{bounded structural relationship}, recur across \result{family recurrence} independent datasets and collectively show \result{family-level effect}, whereas individual features exhibit \result{feature-level heterogeneity}. We use programme to describe this reproducible family-level biological pattern, not to assert a biochemical pathway.

We selected \result{anchor count} commercially available structures using a criterion frozen before purchase. Authentic standards reproduced precursor mass, tandem spectrum and retention/coelution for \result{successful anchor count}, while \result{failed/ambiguous anchor count} remained unconfirmed and stayed in the denominator. One representative family member remained deliberately unresolved: \result{surviving isomer count} alternatives survived all available evidence, and the atlas records \result{predicted discriminating measurement} as the experiment predicted to distinguish them rather than assigning a unique identity.

The global paired analysis and the replicated programme therefore connect the methodological contribution back to biology: inference through a fragmentation world does not merely increase annotation depth, but changes the molecular units available for reproducible metabolomic discovery.

\section*{Discussion}

Mass-spectral structure elucidation is usually framed as a mapping problem: given an observed spectrum, predict the correct molecular structure. Our results support a broader formulation. The spectrum is an incomplete observation of a latent gas-phase reaction world, and prediction, explanation and inverse structure inference can therefore be treated as conditional queries to the same learned chemical process. This shared World supplies an intermediate causal object between molecular structure and spectral observation that direct mappings do not explicitly represent.

The benefit of this formulation is not that every learned trajectory becomes a proven mechanism. Hard atom, electron, charge, hydrogen, mass and valence constraints establish chemical bookkeeping and executability, whereas event propensity and mechanistic relevance remain learned and empirically tested. The value of the World is instead that forward and inverse claims can be required to agree within the same state and transition semantics. A generated molecule is no longer only a high-scoring answer; it is a hypothesis that must produce a reaction network capable of explaining the observation.

This changes how uncertainty is represented. Ambiguous spectra naturally induce multiple molecular hypotheses, and multiple fragmentation paths may support the same observed mass. GFlowNet and sequential posterior inference therefore preserve distributions over trajectories and structures rather than optimizing only one winner. Guided inference further separates chemical probability from evidential value: common fragmentation can be likely yet uninformative, whereas a rarer candidate-specific path can be decisive for structural discrimination. Reinforcement learning is used here to allocate finite computation across this combinatorial world, not to relax chemical constraints.

The resulting evidence hierarchy provides the bridge from method to science. Dark spectra are not all promoted to a rank-1 molecular identity. Instead, each spectrum reaches the strongest level supported by its evidence, from formula compatibility through World reachability and bidirectional consistency to calibrated bounded structure and, for a small subset, authentic-standard confirmation. This makes repository-scale analysis possible without conflating model confidence with experimental identification.

Applied at scale, this approach reveals a recurrent dark chemical population with non-random structural organization and reproducible biological dependence. The distinction between near-known chemistry and remote recurrent families quantifies where existing metabolite references are informative and where they remain sparse. The feature-versus-family analysis further shows why structural resolution matters downstream: chemically related observations can be compared as a common molecular object across cohorts rather than as unrelated anonymous peaks.

Important limitations remain. A learned World is an approximation to gas-phase chemistry; legal electron-flow events are not automatically validated mechanisms. The experimental observation model is incomplete, and unobserved generated states cannot generally be labelled false. Prospective library growth is not a random sample of dark chemistry. Repository-scale entity counts depend on imperfect collapsing rules. Perturbation dependence does not prove biosynthetic origin, and family-level associations do not establish pathways. Authentic standards remain the highest evidence tier for the structures tested here.

The framework nevertheless changes the role of mass-spectral AI. Rather than asking only whether a model can guess a molecule from a spectrum, it asks whether a shared chemical world can support forward prediction, conditional explanation and inverse inference with explicit evidence. That capability turns the dark metabolome from a collection of unidentified spectra into a graded, testable and biologically analysable chemical space.

\section*{Methods}

\subsection*{Study design and evidentiary separation}
\methodtodo{Preregistered confirmatory endpoints; immutable separation of model/program development, calibration, prospective evaluation, census construction and biological analysis; protected-data firewalls preventing prospective structures, source labels or phenotype labels from influencing structural-family construction.}

\subsection*{Spectral corpora, public libraries and frozen dark-subset definition}
\methodtodo{Repository accessions, mouse cohort provenance, GNPS/MassBank/MoNA/NIST versions and snapshot dates; spectrum preprocessing and QC; library-search engine/configuration; false-accepted-match validation; definition of dark as absence of an accepted match under the frozen protocol.}

\subsection*{Leakage control}
\methodtodo{Exact-structure and close-analogue overlap between all evaluation chemistry and molecular/spectral pretraining corpora, including simulator-generated data; structural standardization; MCES/Tanimoto criteria and treatment of salts/stereochemistry.}

\subsection*{A conservation-preserving molecular fragmentation World}
For a precursor molecule \(M\) under acquisition condition \(c\), ORBIT-MS represents gas-phase fragmentation as a latent stochastic reaction world \(\mathcal{W}(M,c)\), rather than identifying the observed spectrum with the reaction process itself. A World state is a closed bond--electron representation containing atom identities, explicit hydrogens, atomic numbers, formal charge and radical state, a symmetric bond--electron matrix, exact monoisotopic masses, charged and neutral connected components, the current ionization/adduct state and acquisition metadata. Neutral coproducts remain part of the closed chemical state even when they are not directly observed. The experimental spectrum is treated as a partial observation generated from this latent reaction world; absence of a peak is therefore not used as a generic negative label for an otherwise legal World state.

Elementary transformations are generated autoregressively as single-electron events rather than selected from a catalogue of fragmentation templates. At each event, the policy selects an electron source from an atom or currently occupied bond, chooses whether the target is an atom or a bond, predicts the target atom(s), and transfers one electron. It then chooses whether to stop the elementary event, repeat the preceding electron transfer, or generate a new electron event. Two repeated single-electron transfers in the same direction recover an ordinary two-electron pair transfer, whereas distinct electron destinations permit homolytic and multi-electron rearrangement patterns to emerge from the same action language. The accumulated electron events are applied simultaneously to the parent bond--electron state by a deterministic compiler. Atom identity, total valence-electron count, charge and bond-electron symmetry are conserved exactly; products that cannot be materialized under the supported valence model are rejected. The maximum number of electron events used during sampling is an engineering compute guard that scales with molecular size and is not interpreted as a chemical reaction-depth limit.

The World is expanded as a multi-branch probabilistic reaction network. Policy STOP is represented as one terminal sibling branch and does not delete competing reaction branches. Distinct electron-event trajectories may reconverge to the same canonical chemical state; their path probabilities are retained separately during execution and summed when constructing the canonical reaction DAG. Frontier scheduling determines which unresolved branch receives the next unit of computation, not which chemistry is allowed to exist.

Crucially, probability mass is not renormalized only over successful materialized products. At each sampled state, the policy mass is partitioned into
\[
P_{\mathrm{stop}}+
P_{\mathrm{valid}}+
P_{\mathrm{invalid}}+
P_{\mathrm{censored}}+
P_{\mathrm{unsampled}}=1.
\]
\(P_{\mathrm{invalid}}\) denotes sampled electron-event trajectories rejected by the chemistry compiler or materializer; \(P_{\mathrm{censored}}\) denotes trajectories that attempted to continue beyond an event or execution budget; and \(P_{\mathrm{unsampled}}\) denotes policy mass not represented by the finite Monte Carlo sample. Invalid mass is treated as failed chemistry, whereas censored and unsampled mass remain unresolved. They are never reassigned to successful children. World probability estimates are therefore reported together with halted, invalid, censored and unresolved probability mass, allowing chemical failure, computational incompleteness and learned termination to remain distinct.

\subsection*{Free, Guided and Inverse inference over one World}
The three WGV modes share the same chemical state and transition semantics but answer different conditional questions.

\paragraph{Free World inference.}
Free inference estimates the fragmentation distribution of a molecule without reading the target spectrum at inference time. The policy generates autoregressive electron-event transitions and branch-local STOP probabilities, and Subtrajectory Balance is used to train probability flow across multiple compatible trajectories rather than optimizing a single best path. Training uses molecule--spectrum aligned data as observational evidence, including observed fragment masses, intensities and ordered parent--child transitions, but the target spectrum is not an input to the Free policy. Generated but unobserved legal states remain unlabeled unless a separate observation model supports a negative claim.

\paragraph{Guided spectrum-conditioned inference.}
Guided inference conditions the same World policy on the observed spectrum. The current production verifier uses graph--spectrum conditioning and a learned guided value to prioritize unresolved frontier states and electron-event proposals under a fixed expansion budget; it does not alter the chemistry compiler or the set of legal transitions. An auxiliary goal-conditioned IQL objective is trained on sampled Guided transitions, but the present production rollout is controlled by the spectrum-conditioned graph policy/value rather than by a separately loaded IQL controller. DAgger and hindsight goal relabelling are therefore not claimed as active components of the current production inference path. Free and Guided are compared under matched chemistry and matched computational budgets so that any gain can be attributed to observation-conditioned allocation of computation.

\paragraph{Inverse molecular inference.}
Inverse inference constructs molecular hypotheses from the precursor formula/mass constraints and the observed spectrum. A factorized graph neural policy predicts construction action type, element, formal charge, radical state, attachment or atom-pair pointers and bond order without enumerating the complete graph-edit action space. Formula inventory, charge-aware valence, exact hydrogen closure, connectivity and materializability are enforced by hard masks and the construction executor. Canonical partial-graph identities merge equivalent states reached through different construction histories. During training and proposal shaping, ordered fragmentation evidence can be represented as leaf--parent--precursor progress so that a partial molecular graph is rewarded for continuously explaining an observed fragmentation hierarchy rather than for matching isolated peak masses. Sequential Monte Carlo maintains multiple candidate modes, and every completed candidate is finally replayed through the forward fragmentation World to obtain bidirectional spectrum--structure evidence.

\subsection*{World-based structural evidence hierarchy}
Structural inference is reported as a hierarchy of evidence rather than as a single candidate score. The frozen levels are: precursor mass or formula compatibility; chemically valid molecular structure; reachability in the learned fragmentation World; executable paths supporting observed fragment evidence; candidate-specific discrimination against structural alternatives; bidirectional spectrum--structure consistency after forward World replay; a calibrated bounded structural statement; and authentic-standard confirmation. Generator recall is reported separately from evidence quality so that failure to propose the true structure is not counted as a World-verification error.

World evidence additionally records how completely the relevant reaction space was explored. For a candidate, supported observations are interpreted together with invalid, event-budget-censored and unsampled policy probability mass. A candidate supported only by a small explored fraction of its policy mass is not treated as equivalent to one for which the compatible reaction world is extensively resolved. Evidence thresholds and any summary statistic using unresolved probability mass are frozen on known-structure development/calibration data before repository-scale dark-spectrum analysis.

\subsection*{Soundness--resolution operating point}
\methodtodo{Calibration cohort; stringency sweep; false-exclusion endpoint; confidence mapping; operating-point selection before atlas analysis; bootstrap/clustered uncertainty; sensitivity to mass tolerance, peak threshold and acquisition metadata.}

\subsection*{Prospective library-growth validation}
\methodtodo{Historical freeze reconstruction; artifact hashes; definition of spectra dark at freeze and structures first deposited later; analogue exclusion; depositing-group independence; outcome revelation procedure; calibration/containment metrics and structural-distance stratification.}

\subsection*{Comparator evaluation}
\methodtodo{Exact-mass, spectrum-similarity, CFM-ID, SIRIUS/CSI:FingerID and other executable comparators; identical spectra/candidate pools where applicable; independent calibration mapping for each comparator; same-formula and near-isomer strata.}

\subsection*{Structural-entity collapsing and census uncertainty}
\methodtodo{Rules for adducts, isotopologues, charge states, in-source fragments and redundant acquisitions; validation on known compounds; over-/under-merging rates; propagation into entity-count uncertainty; entity definition at each structural-resolution level.}

\subsection*{Recurrence and dataset independence}
\methodtodo{Definition of independent datasets; handling of shared samples/laboratories/re-depositions; recurrence threshold; sensitivity analyses; structural-family construction from chemistry only.}

\subsection*{Dark chemical landscape}
\methodtodo{Frozen characterized-metabolite reference set; structural-distance metrics; near-known neighbourhood threshold; transformation vocabulary; treatment of class-level-only statements; family clustering/graph construction; reference-set and metric sensitivity.}

\subsection*{Biological-source attribution}
\methodtodo{Public germ-free/gnotobiotic, antibiotic and dietary-intervention datasets; manual design verification; harmonization; biological sample as unit; batch/platform/study effects; positive-control panel; label permutation; mixed/unresolved rules.}

\subsection*{Feature-versus-family biological analysis}
\methodtodo{Family freezing and hashing before phenotype testing; identical samples/covariates for feature and family arms; preregistered primary reproducibility endpoint; multiplicity control; independent replication; leave-one-cohort-out and metadata-permutation controls.}

\subsection*{Authentic-standard validation}
\methodtodo{Frozen anchor-selection rule; commercial standard sourcing; precursor, MS/MS and retention/coelution criteria; handling of stereoisomers; reporting of failed/ambiguous anchors and alternative candidates.}

\subsection*{Automated World/program evolution}
\methodtodo{Outer-loop failure localization; permitted typed modifications; deterministic chemistry compiler; fresh-data acceptance gates; causal replay, ablation, transfer and regression tests; strict separation between proposal generation and acceptance. This analysis is supporting and reported primarily in Extended Data.}

\subsection*{Statistics}
Unless otherwise stated, the molecule group is the independent unit for structural benchmark analyses, the independent dataset is the unit for recurrence across repositories, and the biological sample is the unit for biological inference. Effect sizes are reported with \result{confidence interval convention}. Paired structural comparisons use \result{bootstrap/sign-flip procedure}. Atlas entity counts propagate \result{collapsing-error uncertainty model}. Biological analyses use \methodtodo{final preregistered multiplicity-control procedure and replication criterion}.

\subsection*{Data availability}
\methodtodo{Repository accessions; mouse cohort deposition/permissions; historical library snapshots; released atlas tables containing structural statement, confidence, recurrence and evidence provenance; public perturbation datasets.}

\subsection*{Code availability}
ORBIT-MS is developed at \texttt{https://github.com/wangyu-sd/orbit-ms}. A frozen release containing model checkpoints, fragmentation-World configuration, candidate-generation configuration, calibration mapping, structural-family definitions, evaluation code and analysis manifests will be archived at \result{Zenodo DOI/release tag}.

\section*{Acknowledgements}
\result{Funding acknowledgements and collaborators.}

\section*{Funding}
\result{Grant numbers and funding agencies.}

\section*{Author contributions}
\result{Author contribution statement.}

\section*{Competing interests}
The authors declare \result{final competing-interests statement}.

\section*{Additional information}
Correspondence and requests for materials should be addressed to \result{corresponding author}.

\bibliographystyle{unsrtnat}
\bibliography{references}

\section*{Extended Data figure legends}

\noindent\textbf{Extended Data Figure 1 | Corpora, splits, leakage and observation contract.}
Dataset provenance; molecule-disjoint splits; historical library snapshots; exact and analogue leakage audits; acquisition conditions; dark-subset definition; distinction between unobserved and negative evidence.

\noindent\textbf{Extended Data Figure 2 | Fragmentation World state, actions and conservation.}
Bond-electron state; explicit atom and H identities; exact mass; charged and neutral products; ionization microstates; sparse electron-flow events; hard action masks; materialization; forward/inverse round trips; action validity and compute scaling.

\noindent\textbf{Extended Data Figure 3 | Free probabilistic reaction-world learning.}
Subtrajectory-balance/GFlowNet training; multibranch siblings; branch-local STOP; reconvergent flow; probability conservation; calibration; observed mass/intensity/transition coverage; direct-spectrum and single-path baselines.

\noindent\textbf{Extended Data Figure 4 | Guided spectrum-conditioned inference.}
Free-versus-Guided matched-budget comparisons; spectrum-conditioned graph/value policy; auxiliary IQL training diagnostics; frontier representation; explanatory path coverage; coverage-vs-compute; target-spectrum leakage firewall; explicit statement that DAgger/HER are not part of the current production rollout.

\noindent\textbf{Extended Data Figure 5 | Inverse molecular inference.}
Heavy-atom construction; canonical state merging; charge/formula/H/valence constraints; GFlowNet/SMC; Exact/Recall@K; structural similarity; posterior diversity; ablations; forward World replay.

\noindent\textbf{Extended Data Figure 6 | Evidence hierarchy and prospective calibration.}
GT/hard-decoy bidirectional discrimination; reliability by evidence level; soundness--resolution trade-off; historical library-growth validation; structural-distance stratification; calibrated comparators.

\noindent\textbf{Extended Data Figure 7 | Outer-loop scientific program evolution.}
Failure localization; typed program/state/action/algorithm revisions; conservation compilation; accepted and rejected proposals; causal replay; paired held-out gates; scientific KPI changes and compute cost.

\noindent\textbf{Extended Data Figure 8 | Repository-scale dark-resolution and entity collapsing.}
Complete evidence-level funnel; unresolved/censored fractions; adduct/isotope/charge/in-source collapse validation; entity-count uncertainty; recurrence and evidence-by-recurrence analyses.

\noindent\textbf{Extended Data Figure 9 | Structural landscape and biological-source organization.}
Known-metabolite reference sensitivity; near-known/remote thresholds; family definitions; public perturbation inventory; positive controls; permutation nulls; source dependence and structure--source coupling.

\noindent\textbf{Extended Data Figure 10 | Feature-versus-family biological analysis.}
Frozen family definitions and contrasts; complete paired feature/family results; cross-cohort replication; effect-sign concordance; multiplicity control; leave-one-cohort-out; metadata permutation and threshold sensitivity.

\noindent\textbf{Extended Data Figure 11 | Authentic-standard and unresolved-case dossiers.}
All selected anchors including failures; candidate sets; World evidence; sample and standard spectra; retention/coelution; surviving alternatives; predicted discriminating measurements.

\end{document}
